\documentclass[10pt,a4paper]{amsart}
\usepackage[margin=19mm]{geometry}
\usepackage{amsmath,amssymb,mathtools}
\usepackage[scr=rsfs,cal=euler]{mathalfa}
\usepackage{xfrac}
\usepackage[unicode,hidelinks,bookmarks=false]{hyperref}
\newcommand{\ie}{i.e.\ }
\newcommand{\cf}{cf.\ }
\newcommand{\resp}{respectively\ }
\newcommand{\Subsection}{\subsection}
\providecommand{\nobreakdash}{\nobreak\hbox{-}}
\setlength{\emergencystretch}{2em}
\multlinegap0pt
\usepackage{float}
\usepackage{amssymb}
\usepackage{bm}
\usepackage{xcolor}
\usepackage{mathtools}
\newtheorem{theorem}{Theorem}
\newtheorem{lemma}{Lemma}
\newcommand{\C}{\mathbb{C}}
\newcommand{\N}{\mathbb{N}}
\newcommand{\Q}{\mathbb{Q}}
\newcommand{\R}{\mathbb{R}}
\title{Connected components of the ranges of twisted divisor functions on number fields}
\author{Sophie Zhu}
\date{Source: arXiv:2605.26482v1 (CC BY 4.0)}
\begin{document}
\maketitle

\noindent\emph{Source and licence.} Connected components of the ranges of twisted divisor functions on number fields. Version arXiv:2605.26482v1 (CC BY 4.0), distributed by arXiv under the Creative Commons Attribution 4.0 International licence, http://creativecommons.org/licenses/by/4.0/, as recorded in the arXiv licence metadata for this exact version and shown on the abstract page https://arxiv.org/abs/2605.26482v1. Full licence notice: \texttt{licenses/arxiv-2605.26482-CC-BY-4.0.txt}.\par\smallskip

\noindent\textit{Editorial scope.} Counted unit: the article's theorem lowerbound1 (source0.tex lines 208--210). The article's complete original proof of it is delivered with it (source0.tex lines 674--698, one \texttt{proof} environment of 25 lines, which the article places away from the statement). No mathematics is written, repaired or re-derived: the statement and the proof are the article's own lines, in the article's own notation, and no retained line is substituted or re-flowed.  Retained uncounted as the source-local context the proof needs: lines 661--667.  Nothing was omitted from any retained span, so the package carries no omission record.  No retained line contains a citation, so the wrapper carries no bibliography and no bibliography entry.  No retained line contains a figure environment, an image-inclusion command or drawing code, so no image is delivered and the package declares no asset dependency.  Editorial formatting only, in addition: an AMS \texttt{amsart} preamble that restates the article's own declarations and macros used by the retained lines, namely the environments \texttt{theorem}, \texttt{lemma} on separate counters, together with the standard packages those lines need.  The retained environments print in this document as Theorem 1, Lemma 1 in order of appearance, whereas the article prints the same items with the numbering of its own full document.  The complete native source of the article as distributed in the arXiv e-print for this exact version is bundled unchanged as \texttt{sources/201419/source0.tex}.
\begin{lemma} \label{tail}
    Let $r > 1.$
    \begin{itemize}
        \item[(1)] For any integer $0 < d \leq r-1,$ we have $$d^{-r} > \sum_{n = d+1}^\infty n^{-r}.$$
        \item[(2)] For any odd integer $d \leq 2r-2,$ we have $$d^{-r} > \sum_{n \geq d+2; 2 \nmid n} n^{-r}.$$ 
    \end{itemize}
\end{lemma}

\begin{theorem} \label{lowerbound1}
    Let $r \geq 3,$ and let $\chi$ be a character of modulus $m$. Then $C_{r,\Q,\chi} \geq 2^{\pi_m(2r-2)-\delta_m},$ where $\delta_m$ equals 0 if $m$ is even and equals 1 otherwise. 
\end{theorem}

\begin{proof}[Proof of Theorem \ref{lowerbound1}]
Let $\theta$ be a nonzero complex number. Then any complex number $x \in \C$ can be written uniquely in the form $a \theta + b \theta i,$ where $a,b \in \R.$ Define $\Re_\theta$ by $\Re_\theta(a\theta + b \theta i) = a.$


We begin by partitioning $\N$. To do so, we first define the map $B_1: \N \rightarrow \mathcal{P}(\N)$ by $B_1(x) = \{d \in \N : d \mid x, d \leq r-1, \gcd(d,m)=1\}.$ Next, we define the map $B_2 : \N \rightarrow \mathcal {P} (\N)$ by $$B_2(n) = \begin{cases}
    \{d \in \N : d \mid n, r-1 \leq d \leq 2r-2, \gcd(d, 2m) = 1\} & \text{ if } 2 \not\in B_1(n) \\
    \emptyset &\text{ if }2 \in B_1(n).
\end{cases}$$
Finally, define $A(n) = (B_1(n), B_2(n)).$
Note that this determines a partition of $\N$ whose parts are $A^{-1}(S)$ for $S \in A(\N).$ We wish to show that for any distinct $S_1, S_2 \in A(\N),$ the sets $\sigma_{r,\Q,\chi}(A^{-1}(S_1))$ and $\sigma_{r,\Q,\chi}(A^{-1}(S_2))$ have disjoint closures. Write $S_1 = A(x)$ and $S_2 = A(y)$ for distinct $x,y \in \N.$ 

    We consider separately the cases $B_1(x) \neq B_1(y)$ and $B_1(x) = B_1(y)$. Assume first that $B_1(x) \neq B_1(y)$.  Suppose without loss of generality that the smallest integer $d_0 \leq r-1$ coprime to $m$ that divides precisely one of $x$ and $y$ lies in $B_1(x)$. Let $\theta = \chi(d_0).$ Take $x'$ such that $B_1(x') = B_1(x)$ and $y'$ such that $B_1(y') = B_1(y).$ For each $d < d_0,$ the difference of the terms for $d$ in $\sigma_{r,\Q,\chi} (x') - \sigma_{r, \Q,\chi}(y')$ is 0: if $\gcd(d,m) > 1,$ then $\chi(d) = 0,$ and otherwise, $d \in B_1(x')$ if and only if $d \in B_1(y')$ by the minimality of $d_0.$ Then we have $$\sigma_{r,\Q,\chi} (x') - \sigma_{r, \Q,\chi}(y') = \theta d_0^{-r} + \sum_{d \mid x'; d > d_0} \chi(d)  d^{-r} - \sum_{d \mid y'; d > d_0} \chi(d) d^{-r}.$$ Applying $\Re_\theta$ to both sides, we obtain $$\Re_\theta(\sigma_{r,\Q,\chi}(x')) - \Re_\theta(\sigma_{r,\Q,\chi}(y')) \geq d_0^{-r} - \sum_{d > d_0} d^{-r} > 0$$ by Lemma \ref{tail}. 
Thus, the closures of $\sigma_{r,\Q,\chi}(A^{-1}(A(x))$ and $\sigma_{r,\Q,\chi}(A^{-1}(A(y))$ are disjoint. 

    It remains to consider the case $B_1(x) = B_1(y).$ Then $B_2(x) \neq B_2(y),$ so one of these sets is nonempty, meaning $2 \not\in B_1(x) = B_1(y).$ If $m$ is even, then $\chi(d) = 0$ for even $d$. If $m$ is odd, because  $r-1 \geq 2$, the fact that $2 \not\in B_1(x) = B_1(y)$ implies that $x$ and $y$ have no even divisors.  The argument then proceeds in a similar fashion to the case $B_1(x) \neq B_1(y).$ Let $d_0$ be the smallest integer that is in precisely one of the two sets $B_2(x)$ and $B_2(y).$ Without loss of generality, we assume $d_0 \in B_2(x) \setminus B_2(y).$ Let $\theta = \chi(d_0),$ and take $x'$ such that $A(x') = A(x)$ and $y'$ such that $A(y') = A(y).$ When comparing $\Re_\theta(\sigma_{r,\Q,\chi}(x'))$ and $\Re_\theta(\sigma_{r,\Q,\chi}(y'))$, we use the fact that the contributions from divisors $d < d_0$ are the same. In addition, there are no contributions from even divisors. Indeed, if $m$ is even, then $\chi(d) = 0$ for even $d$; if $m$ is odd, because  $r-1 \geq 2$, the fact that $2 \not\in B_1(x) = B_1(y)$ implies that $x$ and $y$ have no even divisors.  More explicitly, we compute 
    \begin{align*}
        \Re_\theta(\sigma_{r,\Q,\chi}(x'))-\Re_\theta(\sigma_{r,\Q,\chi}(y')) &= \Re_\theta \left( \sum_{d \mid x'; d < d_0} d^{-r} \chi(d) + d_0^{-r} \theta + \sum_{d \mid x'; d \geq d_0+2; d \text{ odd}} d^{-r}\chi(d)\right) \\
        &\text{ \; } - \Re_\theta \left( \sum_{d \mid y'; d < d_0} d^{-r} \chi(d) +  \sum_{d \mid y'; d \geq d_0+2; d \text{ odd}} d^{-r}\chi(d)\right) \\
        &= d_0^{-r} + \Re_\theta\left( \sum_{d \mid x'; d \geq d_0+2; d \text{ odd}} d^{-r}\chi(d) - \sum_{d \mid y'; d \geq d_0+2; d \text{ odd}} d^{-r}\chi(d)\right) \\
        &\geq d_0^{-r} - \sum_{d \geq d_0 + 2 ; d \text{ odd}} d^{-r},
    \end{align*}
    which is positive by Lemma \ref{tail}. Thus, we have $C_{r, \Q, \chi} \geq |A(\mathbb N)|.$ 

    It remains to give a lower bound for $|A(\N)|.$ Let $P$ be the set of odd primes $p \leq 2r-2$ that are coprime to $m$. For each subset $S \subseteq P,$ define $$n_S = \prod_{p \in S} p .$$ Because $n_S$ is odd, $2 \not\in B_1(n_S),$ so $B_2(n_S) = \{d \in \N : d \mid n_S, r-1 \leq d \leq 2r-2, \gcd(d, 2m) = 1\}.$ Each prime $p \in S$ is recorded either in $B_1(n_S)$ (if $p \leq r-1$) or in $B_2(n_S)$ (if $r-1\leq p \leq 2r-2$). Thus, distinct choices of subsets $S$ of $P$ give rise to distinct values of $A(n_S)$. Consequently, we have $$C_{r, \Q, \chi} \geq |A(\mathbb N)| \geq 2^{|P|} = 2^{\pi_m(2r-2) - \delta_m},$$ where $\delta_m$ equals 0 if $m$ is even and equals 1 if $m$ is odd.
\end{proof}

\end{document}
